% 02-abbreviations.tex - Abbreviations and Definitions

\chapter{Abbreviations and Definitions}
\label{sec:abbreviations}

\section{Abbreviations}

\begin{longtable}{|l|p{10cm}|}
    \hline
    \textbf{Abbreviation} & \textbf{Definition} \\
    \hline
    \endfirsthead
    \hline
    \textbf{Abbreviation} & \textbf{Definition} \\
    \hline
    \endhead
    \hline
    \endfoot
    DR & Data Recorder - recording subsystem running the DROS software \\
    \hline
    DRSU & Data Recorder Storage Unit - removable storage for DRs \\
    \hline
    FIFO & First In, First Out - queue processing order \\
    \hline
    ICD & Interface Control Document \\
    \hline
    LWA & Long Wavelength Array \\
    \hline
    LWA1 & Long Wavelength Array station 1 \\
    \hline
    LWA-NA & Long Wavelength Array North Arm station \\
    \hline
    LWA-SV & Long Wavelength Array Sevilleta station \\
    \hline
    MCS & Monitor and Control System \\
    \hline
    MIB & Management Information Base - hierarchical data structure for status reporting \\
    \hline
    MJD & Modified Julian Date \\
    \hline
    MPM & Milliseconds Past Midnight \\
    \hline
    SMTP & Simple Mail Transfer Protocol \\
    \hline
    SSH & Secure Shell \\
    \hline
    UCF & LWA User Computing Facility - cluster for reducing LWA data \\
    \hline
    UDP & User Datagram Protocol \\
    \hline
    UTC & Coordinated Universal Time \\
    \hline
    WAL & Write-Ahead Logging - SQLite journaling mode for reliability \\
    \hline
    \O{}MQ & \O{}MQ (ZeroMQ) - high-performance asynchronous messaging library \\
    \hline
\end{longtable}

\section{Definitions}

\begin{description}
    \item[Bandwidth Limit] A configurable limit on the transfer rate for remote copies, specified in MB/s. A value of 0 disables limiting.

    \item[Copy Queue] A persistent, per-\DR{} FIFO queue of file transfer operations backed by SQLite. Items survive system restarts.

    \item[Global Inhibit] A per-\DR{} flag that prevents queue processing regardless of \DR{} state. Set by \cmd{PAU} commands and cleared by \cmd{RES} commands.

    \item[Local Copy] A file transfer where the source and destination are on the same host, performed via \code{rsync} without SSH.

    \item[Purge] The process of deleting source files from a \DR{} after they have been successfully copied to a destination.

    \item[Reference ID] A unique integer identifier assigned to each command packet, generated by the \O{}MQ reference server. IDs range from 1 to 999,999,999 and roll over.

    \item[Remote Copy] A file transfer where the source and destination are on different hosts, performed via \code{rsync} over SSH. Only one remote copy may be active at a time across all \DR{} queues.

\end{description}

\section{Document Conventions}

\begin{itemize}
    \item Command names are shown in \cmd{monospace} font (e.g., \cmd{PNG}, \cmd{SCP})
    \item MIB entry names are shown in \mibentry{monospace} font (e.g., \mibentry{SUMMARY})
    \item Hexadecimal values are prefixed with \hex{} (e.g., \hex{00}, \hex{04})
    \item File paths are shown in \file{monospace} font (e.g., \file{defaults.json})
    \item Code identifiers are shown in \code{monospace} font (e.g., \code{InterruptibleCopy})
\end{itemize}
